%!TEX root = ../../template.tex
%%%%%%%%%%%%%%%%%%%%%%%%%%%%%%%%%%%%%%%%%%%%%%%%%%%%%%%%%%%%%%%%%%%%
%% architectural_tradeoffs_and_limitations.tex
%% Architectural Tradeoffs and Limitations section for Proposed Model and Implementation
%%%%%%%%%%%%%%%%%%%%%%%%%%%%%%%%%%%%%%%%%%%%%%%%%%%%%%%%%%%%%%%%%%%%

\section{Architectural Tradeoffs and Limitations}
\label{sec:ProposedModelArchitecturalTradeoffsAndLimitations}

Every architectural decision comes with tradeoffs between simplicity, flexibility, performance, and operational overhead. This section looks at the complexity costs and limitations of the platform and explains why each decision was made despite the costs.

\subsection{Complexity Tradeoffs}
The architecture prioritizes extensibility, independent scalability, and protocol flexibility over simplicity. This results in more components than a minimal setup would need: two brokers connected by a stateless bridge, nine application services each with its own container, an event-driven pipeline with eventual consistency, and a full edge deployment with its own database and broker. The split between ingestion, metadata discovery, and storage is a clear example of this trade-off. It adds operational complexity, but it also keeps normalization, metadata updates, and database writes independent under load. Each of these decisions is justified in the preceding sections based on the requirements they satisfy. The main cost is operational: more containers to orchestrate, more interactions to debug, and more infrastructure to monitor. This cost is most visible in the single-server test environment, where all components share the same host. In a distributed deployment, each service would run on its own infrastructure, and the overhead would be proportional to the workload each component handles.

\subsection{Limitations}
The platform runs every stateful component as a single instance. Kafka has no replication, making it a single point of failure for all internal messaging. TimescaleDB has no read replicas or failover. The edge supports only a single node with no fleet management. Any of these components failing would take the platform down.

Several design choices also prevent horizontal scaling. Both the cloud rule cache and the edge rule synchronization use full-replacement strategies rather than incremental updates, which would become bottlenecks as rule sets grow. Real-time WebSocket routing state and per-device anomaly detection state are kept in process memory, so they are lost on restart and cannot be shared across instances.

The anomaly detection model supports only one statistical method, assumes a normal distribution, and has no awareness of seasonal patterns. Periodic cycles like day/night temperature changes can trigger false anomalies even when values are normal for that time of day.

The choice of \gls{JWT} for device authentication bounds offline resilience. The edge can validate tokens locally using the public key but cannot issue or refresh them, so devices whose tokens expire during an extended outage cannot re-authenticate until the cloud is reachable. Client certificates would not have this limitation, but they require a separate provisioning and revocation infrastructure that the platform does not implement.
